\begin{abstract}
Ring and interferometer test structures monitor the width and thickness of silicon waveguides across a wafer, but
none reports the sidewall angle, and a tilt that a width fit absorbs mispredicts rib devices by hundreds of
picometres. We ask which test structures make the angle identifiable. For an oxide-clad strip the mode is
mirror-symmetric, so to first order a wall tilt is exactly a width change, and no set of strips can
separate the two. A partially etched rib breaks the symmetry, but a rib from a
separate etch brings its own depth and angle, and one rib width then leaves the strip angle unidentifiable even when
the depth is known. Two widths of the same rib restore identifiability. With rings on several strip widths and two
rib widths, the Cram\'er--Rao bound on the strip angle is \qty{0.39}{\degree} at a \qty{1}{pm} resonance-wavelength
precision, with the etch parameters and material indices unknown. The bound assumes that all structures share one
set of process parameters. It doubles once their widths differ by \qty{0.039}{nm}, so a working test structure must
co-locate its members or measure their differences. Results come from finite-element mode solutions, Fisher
information and nonlinear fits to simulated data.
\end{abstract}

\begin{IEEEkeywords}
Silicon photonics, sidewall angle, test structures, identifiability, Fisher information, Cram\'er--Rao bound, ring
resonators, process monitoring.
\end{IEEEkeywords}
